\begin{table}[htbp]
\centering
\caption{Summary Statistics}
\label{tab:summary}
\begin{threeparttable}
\small
\begin{tabular}{lccccc}
\toprule
Variable & N & Mean & SD & Min & Max \\
\midrule
\textit{Panel A: Referendum outcomes} & & & & & \\
Turnout (\%) & 175,405 & 49.90 & 11.69 & 11.48 & 100.00 \\
Yes vote share (\%) & 175,405 & 46.52 & 18.00 & 0.00 & 100.00 \\
\addlinespace
\textit{Panel B: Earthquake salience (instrument)} & & & & & \\
Earthquake salience (7-day) & 175,405 & 22.61 & 17.83 & 8.32 & 116.34 \\
No. earthquakes M5.0+ (7-day) & 175,405 & 30.52 & 23.21 & 16.00 & 142.00 \\
Max magnitude (7-day) & 175,405 & 6.46 & 0.61 & 5.70 & 8.10 \\
\addlinespace
\textit{Panel C: Municipality characteristics} & & & & & \\
Eligible voters & 175,405 & 2594.89 & 7317.87 & 27.00 & 235638.00 \\
Items on ballot & 175,405 & 3.34 & 1.09 & 1.00 & 5.00 \\
\bottomrule
\end{tabular}
\begin{tablenotes}[flushleft]
\footnotesize
\item \textit{Notes:} Sample: 30 Swiss federal referendum dates (2015--2024), municipality level, French- and German-speaking regions. Earthquake salience is the magnitude-weighted inverse-distance score of M5.0+ earthquakes to language-specific media centroids (7 days pre-vote).
\end{tablenotes}
\end{threeparttable}
\end{table}
